% ============================================================
% Methods: Operon Segmentation from PacBio Long-Read RNA-Seq
% ============================================================

\subsubsection{Isoform-based operon segmentation}
\label{method:pacbio-operon}

Of the 267~k clustered RNA isoforms, 4~k supported by $\geq$50 reads were retained to suppress noise from sequencing artifacts and RNase degradation fragments.
These isoforms were grouped into candidate operons by containment clustering: two isoforms were placed in the same cluster if one was fully contained within the other (within a 10~bp boundary tolerance), implemented via a union-find algorithm with a coordinate sweep line. 
This containment clustering collapses nested isoforms — the internal termination and RNA processing products of one operon — into the enclosing unit.
This produced 313 initial operons (147 on the plus strand, 166 on the minus strand; median length 2{,}126~bp), which were then annotated against the 911 canonical \textit{Syn1} genes, of which 304 operons contained at least one sense gene (104 single-gene, 98 two-gene, and 102 with three or more genes).

Adjacent same-strand operons can represent split fragments of a single operon, arising in two ways: either two operons that overlap (29 pairs, 24 sharing a gene), produced when RNase cleavage leaves a separate 5\textquotesingle\ and 3\textquotesingle\ piece, or two operons separated by a short gap ($\leq$600~bp) that still contains a same-strand gene, a clustering coverage hole where no isoform passed the read threshold.
% (e.g.\ \gene{rpsQ}{0662} inside the ribosomal-protein supercluster).
Rather than resolving these by boundary location, each such candidate junction was tested for direct co-transcription from the read data: a pair was merged only when at least 50 strand-specific PacBio reads spanned the junction window ($\pm$80~bp) and the mean read depth across the junction was at least half that of the flanking regions, and was otherwise kept separate as a genuine terminator/promoter boundary.
Of the 61 candidate junctions (29 overlapping, 32 gene-in-gap), 43 passed and 18 were kept separate; merges were applied pairwise with no transitive chaining (each operon joining at most one merge), producing 37 merged operons.
A final union-find pass then merged any same-strand operons sharing identical sense gene loci, combining one further pair and yielding 275 isoform-derived operons.

After the initial segmentation, 215 of 911 genes (23.6\%) remained uncovered by any isoform-derived operon.
These were rescued through three complementary strategies: (i)~20 groups of consecutive uncovered genes (gap $\leq$500~bp) were recovered by identifying spanning isoforms in the BAM file, (ii)~the two rRNA operons (16S--23S--5S) were added manually, as their extreme secondary structure precludes faithful long-read capture, and (iii)~the remaining 162 uncovered genes with non-zero spanning reads were added as single-gene transcription units.
The final operon map comprised 459 operons (237 from isoform clustering, 37 merged, 1 gene-combined, 22 multi-gene rescue, and 162 single-gene rescue) and achieved complete sense-strand coverage of all 911 annotated \textit{Syn1} genes (91.3\% sense-only, 5.6\% both sense and antisense, 3.1\% antisense-only, 0\% uncovered).

\subsubsection{Locate Transcription Promoter and Terminator Signatures}
\label{method:operon-TSSTTS}

Promoter and terminator signatures were characterized only for canonical operons, tentatively defined as isoform-derived operons (segmentation type \texttt{isoform\_operon}) whose transcription start site (TSS) and transcription termination site (TTS) both fall in intergenic regions, meaning that neither boundary lies inside a same-strand gene body.
This restriction retains only operons whose mapped \fiveprime\ and \threeprime\ ends reflect genuine promoter and terminator positions rather than RNase-truncated or internally nested transcript ends.
The TSS was taken as the operon \texttt{start0} on the plus strand and the \texttt{end0} on the minus strand, with the TTS at the opposite boundary.

For each canonical TSS, the upstream promoter was read off the transcribed strand of the \syn~genome (with circular wraparound).
The $-10$ box was located by IUPAC consensus matching, scanning a 6-mer consensus \texttt{TANAAT} over the window from $-12$ to $-7$ and an extended 9-mer consensus \texttt{TNNTANAAT} over $-15$ to $-7$ relative to the TSS, each allowed a register shift of up to $\pm 2$~bp, and selecting the offset that minimized the number of IUPAC mismatches (ties broken by the smallest shift).
The $-35$ region (from $-37$ to $-31$) carried no fixed consensus and was instead summarized by a position-specific sequence logo, which showed it to be AT-rich.
Information-content sequence logos were also generated for the $-10$ box and the TSS context.

Rho-independent (intrinsic) terminators were predicted across the \syn~genome with TransTermHP~\cite{kingsford_rapid_2007}, each prediction consisting of a stem-loop hairpin followed by a \threeprime\ poly-uridine tail and carrying a confidence score.
A predicted terminator was assigned to a canonical operon as terminating when it fell within 50~bp downstream of the TTS in the transcription direction, or as internal when it lay inside the operon body at least 50~bp from the TTS.
The terminator hairpins were rendered as sequence logos.
